\section{Introduction}

\subsection{Study Rationale}
\draftinstr{Follow the Phase 2 publication section
\texttt{final-protocol/publication/sections/sec-02-introduction.tex} and the Phase 1
model file \texttt{trial-protocol/draft-protocol/sections/sec-02-introduction.tex}.
State the PDAC burden (five-year OS $< 13$ percent \cite{Siegel2025}), the
Whipple fistula-sepsis cascade \cite{Bassi2017}, and the 2025 Dutch nationwide
robotic cohort baseline \cite{DutchCohort2025}; establish equipoise from the Phase 1
predicate \cite{kawchak2026phase1} and frame the randomized 1:1 comparison of Arm A
versus Arm B (modified FOLFIRINOX \cite{Conroy2018}). Source from
\texttt{trial-protocol/inputs/2030-pdac-1min-final-paper/sections/\{introduction,discussion\}.tex}.}

\draftinstr{Render \texttt{trial-phase-2/mermaid/fig-03-counterfactual-scenarios.md}
as a TikZ \texttt{mermaidfig} (this is \textbf{Figure 3}): the four counterfactual
scenarios (A resection-window collapse; B intra-op SMV/PV injury averted by the
no-fly gate; C daraxonrasib restart mistiming; D under-funded under-powered study
versus co-investment behind the firewall) in which withholding the integrated
approach worsens the patient or the evidence.}

\subsection{Background}
\draftinstr{Carry the four Background subsubsections: PDAC epidemiology and the
Whipple fistula-sepsis cascade; daraxonrasib mechanism and clinical program (FDA
Breakthrough Therapy \cite{revbreakthrough}, RASolute 302 \cite{rasolute302}, and
the five author simulation works
\cite{kawchak_2026_20196639,kawchak_2025_17239510,kawchak_2025_17001137,kawchak_2025_16415815,kawchak_2025_15735068});
current state of AI in clinical trials and surgery (FDA AI-device taxonomy
\cite{Singh2025} and the author LLM-trust record
\cite{kawchak2025codegen16,kawchak2026unitree,kawchak2026adapt312,kawchak2026hr9510});
and funding / co-investment success likelihood. Source from
\texttt{trial-protocol/inputs/2030-pdac-1min-final-paper/sections/\{introduction,methods,discussion\}.tex}
and \texttt{trial-protocol/inputs/author\_works.bib}.}

\draftinstr{Render \texttt{trial-phase-2/mermaid/fig-04-physical-ai-concerns.md} as
a TikZ \texttt{mermaidfig} (this is \textbf{Figure 4}): the nine Physical AI
concerns mapped to their mitigations (the original eight device concerns plus the
ninth, financial influence and inequitable access). State each pairing in full in
the full-width \textbf{Table~\ref{tab:concerns}}. Source the concerns from
\texttt{trial-protocol/inputs/21cfr312\_adapt/01\_preamble\_scope\_definitions.tex}.}

\draftinstr{Build the patient-aligned co-investment table
(\textbf{Table~\ref{tab:coinvest}}) mapping each tranche (eight-site network,
Phase 0 compute, robot fleet plus redundancy, Patient Access and Equity Fund,
central review and biomarker core, independent oversight) to the operational lever
it funds and to the capital firewall, governed by H.R. 9510 and disclosed under 21
CFR part 54 \cite{kawchak2026hr9510,ecfr2026part54}. Source the firewall logic from
\texttt{trial-protocol/inputs/auto-bill-02/final-bill/main.tex}.}

\subsection{Risk/Benefit Assessment}

\subsubsection{Known Potential Risks}
\draftinstr{Carry the three risk categories: drug risks (daraxonrasib toxicity at
the RP2D, perioperative healing risk); device and robotic risks (malfunction,
erroneous AI advisory, cyber-physical failure, vascular injury); and operative
risks intrinsic to the Whipple (ISGPS fistula \cite{Bassi2017}, hemorrhage,
delayed gastric emptying, sepsis, conversion). Note that Arm B carries the
operative and modified-FOLFIRINOX risks of standard care
\cite{Conroy2018,DutchCohort2025}.}

\subsubsection{Known Potential Benefits}
\draftinstr{Carry the directed benefits: in-window R0 resection; layered hard
safety limits with a full hash-chained audit trail; optimized daraxonrasib restart
timing (T+7/T+14/T+21); and the Patient Access and Equity Fund participant
benefit.}

\subsubsection{Assessment of Potential Risks and Benefits}
\draftinstr{Render \texttt{trial-phase-2/mermaid/fig-05-risk-benefit.md} as a TikZ
\texttt{mermaidfig} (this is \textbf{Figure 5}): weigh the novel risks against the
benefits for this low-survival population and conclude a favorable balance
consistent with 21 CFR \S312.84. Then render
\texttt{trial-phase-2/mermaid/fig-06-coinvestment-firewall.md} as a TikZ
\texttt{mermaidfig} (this is \textbf{Figure 6}): the co-investment-to-success-likelihood
pathway with the dashed blocked edge showing no funder path to endpoints,
adjudication, or analysis \cite{kawchak2026hr9510,ecfr2026part54}.}
